\documentclass[11pt,a4paper]{article}

% --- Packages ---
\usepackage[margin=.75in]{geometry}
\usepackage[utf8]{inputenc}
\usepackage[T1]{fontenc}
\usepackage{lmodern} % Fix for font not found error
\usepackage{titlesec}
\usepackage{enumitem}
\usepackage{hyperref}
\usepackage{xcolor}
\usepackage{fancyhdr}
\usepackage{comment}
\usepackage{cleveref}
\usepackage[bottom]{footmisc}
\usepackage{graphicx}
\usepackage{longtable}
\usepackage{booktabs}
\usepackage{array}
\usepackage{calc}
\usepackage{etoolbox}

% Standard Pandoc setup
\providecommand{\tightlist}{%
  \setlength{\itemsep}{0pt}\setlength{\parskip}{0pt}}

\crefformat{section}{\S#2#1#3}
\crefformat{subsection}{\S#2#1#3}
\crefformat{subsubsection}{\S#2#1#3}
\crefformat{figure}{#2Figure~#1#3}
\crefformat{footnote}{#2\footnotemark[#1]#3}

\linespread{0.9} % Riduce lo spazio tra le linee

% --- Colors ---
\definecolor{primary}{RGB}{0, 0, 0}
\definecolor{links}{HTML}{0000EE}
\definecolor{blue}{HTML}{26428B}

\hypersetup{
    colorlinks=true,
    linkcolor=links,
    urlcolor=links,
    citecolor=links,
}

% --- Section Formatting ---
\titleformat{\section}
  {\normalfont\Large\bfseries\color{blue}}{\thesection}{1em}{}
\titleformat{\subsection}
{\normalfont\large\bfseries\color{blue}}{\thesubsection}{1em}{}
\titleformat{\subsubsection}
{\normalfont\bfseries\color{blue}}{\thesubsubsection}{1em}{}

% --- Header and Footer ---
\pagestyle{fancy}
\fancyhead[L]{\small Giuseppe Capasso}
\fancyhead[R]{\small intro-to-linux}
\fancyfoot[C]{\thepage}
\renewcommand{\headrulewidth}{0.4pt}

% --- Custom Title Command ---
\newcommand{\proposaltitle}[1]{
    \begin{center}
        \vspace*{0.1cm}
        {\LARGE \bfseries \color{blue} #1} \\[1.5ex]
        {\large \textbf{Giuseppe Capasso}} \\[1ex]
                \small{
            \vspace{0.1cm}
            \textbf{Materia:} intro-to-linux\\
        }
            \end{center}
    \vspace{0.3cm}
    \hrule height 0.5pt
    \vspace{0.5cm}
}

\begin{document}

\proposaltitle{Laboratorio Pratico: Consolidamento Lezione 1}

Laboratorio Pratico: Consolidamento Lezione 1 Obiettivo: Padroneggiare
la navigazione nel filesystem, la manipolazione di file/directory e
l'uso dell'autodocumentazione (man pages). Prerequisiti: Macchina
Virtuale Ubuntu avviata e terminale aperto.

Esercizio 1: Il Navigatore Obiettivo: Prendere confidenza con i percorsi
(path) assoluti e relativi. 1. Apri il terminale. Assicurati di essere
nella tua Home Directory (/home/tuonome). 2. Crea una cartella chiamata
Laboratorio\_Ripasso. 3. Entra nella cartella appena creata. 4. Crea tre
sottocartelle: Documenti, Immagini, Script. 5. Entra in Documenti. 6.
Senza usare cd .. o cd /home\ldots, prova a tornare alla cartella
Laboratorio\_Ripasso usando un solo comando. • Suggerimento: Ricorda
cosa significano i due punti \ldots{} 7. Verifica la tua posizione con
pwd. Verifica: Il comando ls -R (dalla tua home) dovrebbe mostrarti la
struttura ad albero creata.

Esercizio 2: Gestione File e ``File Nascosti'' Obiettivo: Creare,
rinominare e comprendere i file nascosti. 1. Spostati nella cartella
Laboratorio\_Ripasso/Script. 2. Crea un file vuoto chiamato mioprog.sh.
3. Crea un altro file chiamato .configurazione (fai attenzione al punto
iniziale). 4. Esegui il comando ls. Vedi entrambi i file? 5. Trova il
comando (o la ``flag'') giusta per vedere tutti i file, inclusi quelli
nascosti. 6. Rinomina mioprog.sh in backup\_prog.sh.

7. Sposta backup\_prog.sh nella cartella ../Documenti (senza uscire
dalla cartella in cui sei). Domanda per lo studente: Perché il file
.configurazione non si vedeva all'inizio? A cosa servono i file che
iniziano con un punto in Linux?

Esercizio 3: RTFM (Read The Manual) Obiettivo: Imparare a usare man per
scoprire opzioni non spiegate a lezione. Sei un sysadmin e devi creare
una struttura di directory profonda, ma sei pigro e non vuoi fare mkdir
tre volte. 1. Apri il manuale di mkdir digitando man mkdir. 2. Cerca
l'opzione (flag) che permette di creare directory ``parent'' (genitori)
automaticamente se non esistono. • Indizio: Cerca la parola ``parents''.
3. Una volta trovata l'opzione, crea questa struttura con un solo
comando: Laboratorio\_Ripasso/ProgettoA/Sorgenti/Codice (Nota: Le
cartelle ProgettoA e Sorgenti non esistono ancora). 4. Usa il comando
tree (se installato) o ls -R per verificare il successo.

Esercizio 4: Manipolazione Ricorsiva (Attenzione!) Obiettivo: Copiare e
cancellare intere directory. Questo è il punto dove spesso si fanno
errori. 1. Torna nella cartella Laboratorio\_Ripasso. 2. Crea un file di
testo leggimi.txt dentro la cartella ProgettoA (creata nell'esercizio
3). 3. Prova a copiare l'intera cartella ProgettoA dentro Documenti
usando cp ProgettoA Documenti/. • Osserva l'errore: Il terminale ti dirà
che stai omettendo una directory. 4. Usa man cp per trovare l'opzione
``recursive'' (ricorsiva). 5. Riprova il comando aggiungendo la flag
corretta. 6. Ora prova a cancellare la cartella originale ProgettoA con
rm. • Osserva l'errore: Anche qui, rm si rifiuta di cancellare una
cartella piena.

7. Trova l'opzione per cancellare ricorsivamente e rimuovi ProgettoA.
Attenzione: Quando usi rm con le opzioni ricorsive, fai molta attenzione
a dove ti trovi!

Esercizio 5: Il ``Disaster Recovery'' (Challenge Finale) Scenario: Un
collega distratto ha creato dei file importanti ma li ha nominati male e
messi nel posto sbagliato. Devi sistemare tutto. Setup (Esegui questi
comandi per creare il ``disastro''): Bash cd
\textasciitilde/Laboratorio\_Ripasso mkdir -p Disastro/Temp touch
Disastro/Temp/file\_imporante.txt \# (Notare l'errore di battitura)
touch Disastro/Temp/RELAZIONE\_FINALE.doc touch
Disastro/NON\_TOCCARE.txt

La tua Missione: 1. Sposta RELAZIONE\_FINALE.doc nella cartella
Documenti (che hai creato nell'Ex 1). 2. Entra in Disastro/Temp. 3.
Correggi il nome di file\_imporante.txt in file\_importante.txt. 4.
Copia file\_importante.txt nella tua Home directory (\textasciitilde).
5. Torna alla radice di Laboratorio\_Ripasso. 6. Elimina interamente la
cartella Disastro e tutto il suo contenuto. Controllo Finale: Alla fine,
dentro Laboratorio\_Ripasso dovresti avere solo le cartelle Documenti,
Immagini, Script e la struttura interna pulita.



\end{document}
